\section{Optimization by an exact penalty}
\label{d:sec-optimization}

Retain the notation of Section~\ref{d:sec-objective}: \(n\ge4\) is even,
\(P=P(K_n)\), \(p=\binom n2\), and
\(f^*=\max_{x\in P}f(x)\).
We first replace optimization over \(P\) by optimization over the box
\([0,1]^p\).  For \(x\) in this box, define
\begin{equation}
\label{d:eq-max-violation}
w(x)=
\max\left\{
0,\ 
\max_{i\in V}\bigl|x(\delta(i))-1\bigr|,\ 
\max_{\substack{S\subseteq V\\ |S|\ \text{odd}}}
       \bigl(1-x(\delta(S))\bigr)
\right\}.
\end{equation}
By Theorem~\ref{geom:thm-polytope}, \(w(x)=0\) exactly when \(x\in P\).

\begin{lemma}[Feasibility repair]
\label{d:lem-feasibility-repair}
Set \(A_0=10n^2\).  For every \(x\in[0,1]^p\), there is \(y\in P\) such that
\[
\|y-x\|_\infty\le (2+A_0)w(x).
\]
\end{lemma}

\begin{proof}
Write \(w=w(x)\), \(d_i=1-x(\delta(i))\), and
\(d_{\mathrm{tot}}=\sum_{i\in V}d_i\).  Correct the degrees by setting
\begin{equation}
\label{d:eq-degree-correction}
x'_{ij}
=x_{ij}+\frac{d_i+d_j}{n-2}
          -\frac{d_{\mathrm{tot}}}{(n-2)(n-1)}
\qquad (i\ne j).
\end{equation}
The correction sums to \(d_i\) at vertex \(i\), since
\[
\sum_{j\ne i}(d_i+d_j)=(n-2)d_i+d_{\mathrm{tot}}.
\]
Thus \(x'(\delta(i))=1\) for every \(i\).  Since
\(|d_i|\le w\) and \(|d_{\mathrm{tot}}|\le nw\),
\[
\|x'-x\|_\infty
\le
\left(\frac{2}{n-2}+\frac{n}{(n-2)(n-1)}\right)w
\le 2w.
\]
Consequently \(x'_e\ge-2w\), and every odd set \(S\) satisfies
\begin{equation}
\label{d:eq-corrected-cut-bound}
x'(\delta(S))\ge 1-(1+2p)w.
\end{equation}

Let \(u_{ij}=1/(n-1)\).  This vector has degree sums one, and
\[
u(\delta(S))=\frac{|S|(n-|S|)}{n-1}.
\]
Cuts with \(|S|=1\) or \(|S|=n-1\) have value one.  For every nontrivial
odd cut, \(3\le |S|\le n-3\), and hence, when \(n\ge6\),
\[
u(\delta(S))\ge\frac{3(n-3)}{n-1}\ge\frac32.
\]
There are no nontrivial odd cuts when \(n=4\).  In particular, \(u\in P\).

Set
\begin{equation}
\label{d:eq-uniform-repair}
t=\min\{1,A_0w\},
\qquad
y=(1-t)x'+tu.
\end{equation}
If \(A_0w\ge1\), including equality, then \(y=u\in P\).
Otherwise \(t=A_0w<1\).  The inequalities
\[
\frac{A_0}{n-1}\ge2,
\qquad
\frac{A_0}{2}\ge1+2p
\]
give
\[
y_e\ge-2(1-t)w+\frac{A_0w}{n-1}\ge0.
\]
Equation~\eqref{d:eq-corrected-cut-bound} also gives, for each nontrivial
odd cut,
\[
y(\delta(S))
\ge 1-(1-t)(1+2p)w+\frac{A_0w}{2}
\ge1.
\]
The degree sums remain one, so the trivial odd cuts also have value one.
Theorem~\ref{geom:thm-polytope} now yields \(y\in P\) in this branch
as well.

For either branch, use the identity
\[
y-x=(1-t)(x'-x)+t(u-x).
\]
Both \(u\) and \(x\) lie in the box, so
\[
\|y-x\|_\infty
\le 2(1-t)w+t\|u-x\|_\infty
\le 2w+t
\le (2+A_0)w.
\]
This also covers \(w=0\).
\end{proof}

The repair bound converts constraint violation into a quantitative distance
from feasibility. We can therefore penalize violations by a fixed amount
large enough to recover exactly the constrained optimum.

Define the positive integers \(M_0,D_0\) and the function \(\Phi\) by
\begin{equation}
\label{d:eq-penalty-parameters}
M_0=p(K+L)(2+A_0),\qquad
\Phi(x)=f(x)-M_0w(x),\qquad
D_0=p(K+L+M_0)^2.
\end{equation}
A vector \(g\) is a \emph{supergradient} of \(\Phi\) at \(x\) on the box if
\[
\Phi(z)\le \Phi(x)+g\cdot(z-x)
\qquad\text{for every }z\in[0,1]^p.
\]

\begin{lemma}[Exact penalty and integral supergradients]
\label{d:lem-exact-penalty}
The function \(\Phi\) is concave on \([0,1]^p\), and
\[
\max_{x\in[0,1]^p}\Phi(x)=f^*.
\]
At every point of the box, \(\Phi\) has an integral supergradient \(g\)
satisfying \(\|g\|_2^2\le D_0\).
\end{lemma}

\begin{proof}
By Lemma~\ref{d:lem-dyadic-entropy}, the slopes of \(\psi\) are integers of
absolute value at most \(K\).  Since \(|\beta_e|\le L\), it follows that
\begin{equation}
\label{d:eq-objective-lipschitz}
|f(x)-f(y)|\le (K+L)\sum_e|x_e-y_e|
            \le p(K+L)\|x-y\|_\infty
\qquad (x,y\in[0,1]^p).
\end{equation}
For the point \(y\in P\) supplied by Lemma~\ref{d:lem-feasibility-repair},
Equation~\eqref{d:eq-objective-lipschitz} gives
\[
\Phi(x)
\le f(y)+p(K+L)\|x-y\|_\infty-M_0w(x)
\le f(y)\le f^*.
\]
Conversely, \(\Phi=f\) on \(P\), proving equality of the maxima.
The function \(w\) is a maximum of affine functions and is therefore
convex.  Since \(f\) is concave, \(\Phi\) is concave.

Choose a supporting piece slope \(\sigma_e\) of \(\psi\) at each \(x_e\);
at a knot either adjacent slope is allowed.  With
\(b_e=\sigma_e+\beta_e\), these choices give
\[
f(z)\le f(x)+b\cdot(z-x)
\qquad (z\in[0,1]^p).
\]
Write each absolute value in Equation~\eqref{d:eq-max-violation} as the
maximum of its two affine terms.  Choose any active affine term
\(a\cdot z+c\) in this maximum, including the zero term when it is active.
Then \(a_e\in\{-1,0,1\}\), and
\[
w(z)\ge a\cdot z+c=w(x)+a\cdot(z-x).
\]
Subtracting the last inequality, multiplied by \(M_0\), proves that
\(g=b-M_0a\) is a supergradient of \(\Phi\) on the entire box.
The independent choices of supporting pieces and active rows are valid
also at endpoints and ties.  Finally, \(g\) is integral and
\[
|g_e|\le K+L+M_0,
\qquad
\|g\|_2^2\le p(K+L+M_0)^2=D_0.
\]
\end{proof}

We use the classical projected subgradient method
\citep[Equation~(1)]{Polyak1967} to approximate the maximum of \(\Phi\).
Its bounded integral supergradients permit a fixed step size and keep
every iterate on one rational grid. The next lemma gives the finite
iteration count and value guarantee directly.

\begin{lemma}[Finite projected ascent]
\label{d:lem-projected-ascent}
Let \(N_{\mathrm{it}}=pD_0\), set \(x^{(0)}=0\), and for
\(j=0,\ldots,N_{\mathrm{it}}-1\) choose a supergradient
\(g^{(j)}\) as in Lemma~\ref{d:lem-exact-penalty} and set
\begin{equation}
\label{d:eq-projected-iterations}
x^{(j+1)}
=\Pi_{[0,1]^p}\left(x^{(j)}+\frac{g^{(j)}}{D_0}\right),
\end{equation}
where \(\Pi_{[0,1]^p}\) clips each coordinate to \([0,1]\).
Then
\begin{equation}
\label{d:eq-optimization-accuracy}
R=\left\lfloor\max_{0\le j<N_{\mathrm{it}}}\Phi(x^{(j)})\right\rfloor
\qquad\text{satisfies}\qquad
R\le f^*\le R+2.
\end{equation}
\end{lemma}

\begin{proof}
Choose \(z\in P\) with \(f(z)=f^*\); then \(\Phi(z)=f^*\) by
Lemma~\ref{d:lem-exact-penalty}.  Clipping cannot increase distance to
\(z\in[0,1]^p\).  Expanding the squared distance and using the
supergradient inequality therefore gives
\begin{align}
\|x^{(j+1)}-z\|_2^2
&\le \|x^{(j)}-z\|_2^2
 +\frac{2}{D_0}g^{(j)}\cdot(x^{(j)}-z)
 +\frac{\|g^{(j)}\|_2^2}{D_0^2}\notag\\
&\le \|x^{(j)}-z\|_2^2
 -\frac{2}{D_0}\bigl(f^*-\Phi(x^{(j)})\bigr)+\frac1{D_0}.
\label{d:eq-ascent-step}
\end{align}
Summing Equation~\eqref{d:eq-ascent-step} over exactly the indicated
\(N_{\mathrm{it}}\) indices yields
\begin{align*}
2\sum_{j=0}^{N_{\mathrm{it}}-1}
  \bigl(f^*-\Phi(x^{(j)})\bigr)
&\le D_0\bigl(\|x^{(0)}-z\|_2^2
             -\|x^{(N_{\mathrm{it}})}-z\|_2^2\bigr)
       +N_{\mathrm{it}}\\
&\le pD_0+N_{\mathrm{it}}=2N_{\mathrm{it}}.
\end{align*}
Here \(\|x^{(0)}-z\|_2^2\le p\).
Thus the largest of the sampled values, say \(V\), satisfies
\(f^*-1\le V\le f^*\).  Since \(R=\lfloor V\rfloor\),
\[
R\le V\le f^*,
\qquad
f^*\le V+1<R+2,
\]
which proves Equation~\eqref{d:eq-optimization-accuracy}.
\end{proof}

Every iterate in Equation~\eqref{d:eq-projected-iterations} has coordinates
that are integer multiples of \(1/D_0\): the update adds an integral
vector divided by \(D_0\), and clipping inserts only \(0\) or \(1\).
To evaluate \(w\) and choose an active row, compare the zero term and the
\(2n\) signed degree terms with the largest odd-cut violation.
Lemma~\ref{d:lem-odd-cut-oracle} computes a minimum odd cut from the
nonnegative integer capacities \(D_0x_e^{(j)}\) and returns its row in the
original edge coordinates.  The values of \(\Phi\) and the floor defining
\(R\) can then be computed exactly using rational arithmetic.
Proposition~\ref{d:prop-bit-complexity} gives the resulting bit-complexity
bound.
